\section{Gaussian hierarchies and quantitative free energy comparison}
\label{static-sec:gaussian}

The overlap estimates in the later sections require two pressure comparisons:
an upper bound with prescribed replica overlaps, and a lower bound for the
unconstrained pressure with a controlled finite-size error. We derive both
from finite Gaussian hierarchies. The upper bound allows the hierarchy's
covariance endpoint to differ from the prescribed overlap matrix; the cost
of that difference remains explicit. For the lower bound, we choose the
hierarchy along an overlap-dependent differential equation and estimate its
error before making that choice.

The constrained comparison belongs to Guerra's Gaussian interpolation
framework~\cite{Guerra03}, including Talagrand's bounds for coupled
replicas~\cite[Propositions~3.4 and~4.4]{Talagrand07} and Panchenko's
vector-spin formulation~\cite{Panchenko18}. Talagrand's bounds already
retain a cost when prescribed overlaps differ from the covariance
endpoint. We derive the finite-hierarchy matrix form needed here
directly, allowing nonunit endpoint diagonals as well. The quantitative
lower comparison is proved separately below.
We fix $\beta>0$ throughout. A constant $C_\beta$ may vary between
occurrences, but never depends on $n$.

We first remove a constant term from the covariance calculation. Add
independent standard Gaussian variables $J_{11},\ldots,J_{nn}$ to the
disorder and, only in free-energy calculations, use the Hamiltonian
\[
 \widehat H_n(\sigma)
 =H_n(\sigma)+\frac1{\sqrt{2n}}\sum_{i=1}^nJ_{ii}.
\]
The extra term has mean zero and is the same for every $\sigma$.
The Gibbs measure and the expected log partition function are therefore
unchanged. The covariance now takes the exact form
\begin{equation}\label{static-eq:hamiltonian-covariance}
 \E\widehat H_n(\sigma)\widehat H_n(\tau)
 =\frac n2R(\sigma,\tau)^2.
\end{equation}

\subsection{Two Gaussian tools and a hierarchy identity}

We will need concentration for Gaussian disorder and a variance estimate
after conditioning on a spin configuration. The following standard
Gaussian and Brascamp--Lieb inequalities supply these two estimates.

\begin{lemma}[Gaussian tools]\label{static-lem:gaussian-tools}
Let $G$ be a standard Gaussian vector in a finite-dimensional Euclidean space.
If $f$ is $L$-Lipschitz, then for $a>0$ and $\theta\in\R$,
\begin{align}
 \Prob\bigl[|f(G)-\E f(G)|\ge a\bigr]
 &\le 2\exp\left(-\frac{a^2}{2L^2}\right),
 \label{static-eq:gaussian-concentration}\\
 \log\E\exp\{\theta(f(G)-\E f(G))\}
 &\le\frac{\theta^2L^2}{2}.
 \label{static-eq:gaussian-mgf}
\end{align}
If a probability density on $\R^m$ is proportional to $e^{-V}$, where $V$ is
twice differentiable and $\nabla^2V\succeq I$, then
\begin{equation}\label{static-eq:brascamp-lieb}
 \Var(f)\le\E\|\nabla f\|^2.
\end{equation}
In particular, the variance of a linear functional $a\cdot G$ under this
density is at most $\|a\|^2$.
\end{lemma}

The variance bound is the strongly log-concave case of the
Brascamp--Lieb inequality~\cite{BL76}. For clarity, we prove the
concentration estimate and the specialization of the variance bound to
linear functions, which is the only specialization used below.

\begin{proof}[Proof of the needed Gaussian tools]
For the concentration estimate, write $H_s$ for convolution with a
centered Gaussian of covariance $sI$, and let $B_t$ be standard Brownian
motion. The heat equation and It\^o's formula represent a smooth
$L$-Lipschitz function as
\[
 f(B_1)-\E f(B_1)
 =\int_0^1\nabla H_{1-t}f(B_t)\cdot\dd B_t,
 \qquad |\nabla H_s f|\le L.
\]
The integrand is bounded, so its exponential martingale has expectation
one. Its quadratic variation is at most $L^2$; this proves
Equation~\eqref{static-eq:gaussian-mgf} for every real $\theta$.
Apply exponential Markov's inequality and optimize separately for the
two tails to obtain Equation~\eqref{static-eq:gaussian-concentration}.
Smooth convolution and localization, followed by Gaussian
integrability, extend the argument to Lipschitz functions of at most
linear growth. Every application here lies in that class.

For the linear variance estimate, we may work with the potentials that
occur in the proof. They are smooth and satisfy
$I\preceq\nabla^2V\preceq CI$, where the finite constant $C$ may depend
on the fixed finite system. To see the upper bound, note that the
nonquadratic terms are finite logarithmic partition functions or
hierarchy prefix values: their Hessians are bounded sums of
covariances. A linear tilt leaves the Hessian unchanged. Consider the
diffusion
\[
 \dd Y_t=\sqrt2\,\dd B_t-\nabla V(Y_t)\dd t
\]
and denote its semigroup by $Q_t$. The drift is globally Lipschitz,
which gives a unique solution. Under synchronous coupling, strong
convexity contracts distances by $e^{-t}$. Thus an $L$-Lipschitz
function becomes $e^{-t}L$-Lipschitz under $Q_t$. Integration by parts
shows that $\nu(\dd y)\propto e^{-V(y)}\dd y$ is invariant and that
\[
 \frac{\dd}{\dd t}\Var_\nu(Q_tf)
 =-2\int|\nabla Q_tf|^2\dd\nu.
\]
The Gaussian tails supplied by strong convexity justify the same
identity for $f(y)=a\cdot y$ by a cutoff argument. For independent
$Y,Y'$ with law $\nu$, contraction also gives
\[
 \Var_\nu(Q_tf)
 \le\tfrac12 e^{-2t}|a|^2\E|Y-Y'|^2\longrightarrow0.
\]
We can therefore integrate the dissipation identity to infinity.
The gradient bound $|\nabla Q_tf|\le e^{-t}|a|$ yields
$\Var_\nu(a\cdot Y)\le2|a|^2\int_0^\infty e^{-2t}\dd t=|a|^2$.
This proves the linear estimate needed for the conditional disorder
laws. The finite upper Hessian bound justifies the diffusion argument;
at each application we will check its lower Hessian bound explicitly.
\end{proof}


We next define the hierarchy and its probability law. Besides any
quenched root randomness, take independent Gaussian fields
$g_0,\ldots,g_k$. The field $g_0$ belongs to the quenched root. Choose
masses
\[
 0=x_{-1}\le x_0\le\cdots\le x_{k-1}\le x_k=1,
 \qquad w_l=x_l-x_{l-1},
\]
and a terminal log partition function $F_k$ on a finite configuration
space. Integrate the fields backwards by setting, for $1\le j\le k$,
\begin{equation}\label{static-eq:hierarchy-recursion}
 F_{j-1}=\frac1{x_{j-1}}\log\E_j e^{x_{j-1}F_j}.
\end{equation}
Here $\E_j$ integrates the field $g_j$. At a zero mass the right-hand
side is interpreted as $\E_jF_j$. After these recursive integrations,
the hierarchy pressure is $\E F_0$, with expectation also over all
quenched randomness.

The recursion also specifies a probability law on fields and terminal
configurations. Conditional on the prefix through level $j-1$, multiply
the original Gaussian law of $g_j$ by the density
\begin{equation}\label{static-eq:hierarchy-tilt}
 W_j=\exp\{x_{j-1}(F_j-F_{j-1})\}.
\end{equation}
Equation~\eqref{static-eq:hierarchy-recursion} normalizes this density.
After sampling level $k$, draw a configuration from the terminal Gibbs
law. For a pair of paths indexed by a level $l$, sample their root and
fields through $l$ once, and then sample the two continuations
independently conditional on this prefix. We denote expectation under
that pair law by $\E_l^{(2)}$. At $l=k$ the only independent choices
are the terminal thermal draws. More generally, a finite tree is
sampled by drawing each shared prefix once and making its descendants
conditionally independent. In particular, every branch has the
single-path marginal law. This convention fixes the shared randomness
in all subsequent integration-by-parts calculations.

\begin{lemma}[Hierarchy differentiation]\label{static-lem:hierarchy-derivative}
Suppose an independent standard Gaussian vector $g$ at level $j$ enters the
terminal exponent as $\sqrt v\,g\cdot e(S)$, where $e$ is a fixed vector-valued
function of the terminal configuration.  With all other parameters fixed,
\begin{equation}\label{static-eq:hierarchy-derivative}
 \frac{\partial}{\partial v}\E F_0
 =\frac12\left(
 \E\|e(S)\|^2-
 \sum_{l=j}^k w_l\E_l^{(2)}e(S^1)\cdot e(S^2)
 \right).
\end{equation}
A quenched Gaussian is treated as level $0$.  The identity holds for
$v>0$ and extends to interpolation endpoints by integration and continuity.
\end{lemma}

\begin{proof}
To differentiate the pressure, first add a deterministic vector in the
same direction as $\sqrt v\,g$. Differentiating the terminal log
partition function and then the hierarchy recursion shows that, for
$l\ge j$, the gradient of $F_l$ is the conditional mean $m_l^e$ of
$e(S)$ given the prefix through $l$. A second differentiation gives
\begin{equation}\label{static-eq:hessian-recursion-general}
 \nabla^2F_l
 =\E[\nabla^2F_{l+1}\mid l]
   +x_l\Cov(m_{l+1}^e\mid l).
\end{equation}
All conditional expectations in this calculation use the path law.
At the terminal level, the Hessian is the thermal covariance of $e(S)$.

For $j\ge1$, differentiate the $j$th logarithmic expectation with
respect to the variance and integrate by parts in its Gaussian field. The result
is one half of the tilted expectation of
\[
 \tr\nabla^2F_j+x_{j-1}\|m_j^e\|^2.
\]
At level $j=0$, ordinary Gaussian differentiation omits the second
term, exactly as prescribed by $x_{-1}=0$. Earlier levels propagate
the derivative by conditional expectation. Iteration of
Equation~\eqref{static-eq:hessian-recursion-general} now expresses twice
the pressure derivative as
\[
 \E\|e(S)\|^2-\E\|m_k^e\|^2
 +\sum_{l=j}^{k-1}x_l
       \bigl(\E\|m_{l+1}^e\|^2-\E\|m_l^e\|^2\bigr)
 +x_{j-1}\E\|m_j^e\|^2.
\]
The sum telescopes to
$\E\|e(S)\|^2-\sum_{l=j}^kw_l\E\|m_l^e\|^2$.
Two descendants are independent conditional on their common prefix,
so $\E\|m_l^e\|^2=\E_l^{(2)}e(S^1)\cdot e(S^2)$. This is the
claimed identity.

Gaussian integrability justifies each differentiation: the configuration
space and the number of levels are finite, the coupled gradients are
bounded, and the recursive values grow at most linearly in the Gaussian
variables. A zero mass can be handled directly by the expectation
recursion, or obtained as a limit. Integration and continuity give the
stated endpoint interpretation.
\end{proof}

\subsection{A matrix interpolation bound}

We now apply the differentiation identity to a prescribed overlap
constraint. For $d$ configurations $\sigma^1,\ldots,\sigma^d$, write
$D_{ab}=R(\sigma^a,\sigma^b)$ for their Gram matrix. If $D$ is feasible,
meaning that at least one such tuple realizes it, define
\begin{equation}\label{static-eq:constrained-pressure}
 F_{d,n}(D)=\frac1n\E\log
 \sum_{\substack{\sigma^1,\ldots,\sigma^d\in\Sig\\
                  R(\sigma^a,\sigma^b)=D_{ab}\ \forall a,b}}
       \exp\left\{\beta\sum_{a=1}^dH_n(\sigma^a)\right\}.
\end{equation}
The sum is unnormalized. We use the Frobenius notation
$U:V=\sum_{a,b}U_{ab}V_{ab}$ and $|U|^2=U:U$.
To compare this pressure with a single-site recursion, choose a symmetric
matrix path
\[
 Q_{-1}=0,\qquad Q_j-Q_{j-1}\succeq0\quad(0\le j\le k),
 \qquad Q_k=Q_{\rm end}.
\]
Within these conditions, the endpoint is free: it may differ from $D$
and may have nonunit diagonal. Increments of zero covariance are
allowed. If the first nonzero increment should carry positive mass,
insert a quenched level $Q_0=0$; the first genuine increment is then
integrated at mass $x_0$. At a single site, take independent
$d$-dimensional Gaussian fields of covariance
$\beta^2(Q_j-Q_{j-1})$. For a symmetric multiplier $\lambda$, define
$\Psi(Q,x,\lambda)$ as the expected hierarchy value with terminal
function
\[
 \log\sum_{s\in\{-1,1\}^d}
     \exp\left\{\left(\sum_{j=0}^kg_j\right)\cdot s
                    +s^{\mathsf T}\lambda s\right\}.
\]
For the deterministic correction to the pressure, set
\begin{equation}\label{static-eq:matrix-path-integral}
 \int x\,\dd|Q|^2
 =\sum_{j=0}^kx_{j-1}\bigl(|Q_j|^2-|Q_{j-1}|^2\bigr).
\end{equation}
The initial quenched increment therefore carries mass zero in this
integral. The next bound keeps the cost of the free endpoint explicit.

\begin{proposition}[Matrix interpolation]\label{static-prop:matrix-bound}
For every feasible $D$ and every path and multiplier as above,
\begin{equation}\label{static-eq:matrix-bound}
 F_{d,n}(D)\le
 \Psi(Q,x,\lambda)-\lambda:D
 +\frac{\beta^2}{4}
       \left(|D-Q_{\rm end}|^2-\int x\,\dd|Q|^2\right).
\end{equation}
For fixed $D$ and $\lambda$, the bound passes to any sequence of finite
paths for which $Q_{\rm end}$, $\Psi(Q,x,\lambda)$, and the displayed
path integral converge.
\end{proposition}

\begin{proof}
Place independent copies of the single-site fields at all $n$ sites.
Interpolate from those fields to the common Hamiltonian disorder,
retaining the overlap constraint throughout. The terminal exponent is
\[
 \sqrt t\,\beta\sum_{a=1}^d\widehat H_n(\sigma^a)
 +\sqrt{1-t}\sum_{i=1}^n\sum_{j=0}^kg_{j,i}\cdot s_i
 +\sum_{i=1}^ns_i^{\mathsf T}\lambda s_i,
 \qquad s_i=(\sigma_i^1,\ldots,\sigma_i^d).
\]
Denote the expected recursive value, divided by $n$, by $\varphi(t)$.
For two terminal draws, define the cross-overlap matrix by
$R_{ab}=R(\sigma^a,\tau^b)$; this matrix need not be symmetric.
Equation~\eqref{static-eq:hamiltonian-covariance} gives Hamiltonian
covariance $n\beta^2|R|^2/2$, whereas the $j$th site-field increment
has covariance $n\beta^2(Q_j-Q_{j-1}):R$. The self-overlap is $D$.
Sum Lemma~\ref{static-lem:hierarchy-derivative} over these Gaussian
increments to obtain
\begin{align*}
 \varphi'(t)
 &=\frac{\beta^2}{4}\left[
 |D|^2-2Q_k:D
 -\sum_{l=0}^kw_l\E_l^{(2)}
       \bigl(|R|^2-2Q_l:R\bigr)\right]\\
 &\le\frac{\beta^2}{4}
       \left(|D|^2-2Q_k:D+\sum_{l=0}^kw_l|Q_l|^2\right).
\end{align*}
The inequality discards the nonnegative squares $|R-Q_l|^2$ and
therefore does not require symmetry of $R$. Summation by parts gives
\[
 |Q_k|^2-\sum_{l=0}^kw_l|Q_l|^2=\int x\,\dd|Q|^2.
\]
It remains to identify the endpoints. At $t=1$ the fields vanish and
$\varphi(1)=F_{d,n}(D)+\lambda:D$. At $t=0$, remove the overlap
constraint. This increases the terminal sum and hence every recursive
value, because each integration is monotone. The resulting recursion
factors over sites, giving $\varphi(0)\le\Psi(Q,x,\lambda)$.
Integration of the derivative bound proves the proposition. Singular
covariance increments are permitted by continuity, or directly by
representing a field through a square root of its covariance.
\end{proof}

\subsection{The scalar functional}

The scalar version of the hierarchy gives the variational quantity
against which we compare the unconstrained pressure. Let $\mu$ be a
probability measure on $[0,1]$ and put $\alpha(s)=\mu([0,s])$.
Define $\Phi_\alpha(s,z)$ by the terminal condition
$\Phi_\alpha(1,z)=\log(2\cosh z)$ and the backward equation
\begin{equation}\label{static-eq:scalar-pde}
 -\partial_s\Phi_\alpha
 =\frac{\beta^2}{2}
       \left(\partial_{zz}\Phi_\alpha
                    +\alpha(s)(\partial_z\Phi_\alpha)^2\right).
\end{equation}
For a step function $\alpha$, the equation means the Gaussian
recursion already defined: over an interval of length $b-a$ with
constant mass $x$, its terminal function $f$ is mapped to
\[
 z\longmapsto\frac1x\log\E
       \exp\{x f(z+\beta\sqrt{b-a}\,G)\},
\]
with ordinary expectation at $x=0$. For general $\alpha$, take
$L^1$ limits of step functions; Lemma~\ref{static-lem:scalar-regularity}
proves that this definition exists, is independent of the approximation,
and has the stability properties we use. Define
\begin{equation}\label{static-eq:scalar-functional}
 \cP_\beta(\mu)=\Phi_\alpha(0,0)
                   -\frac{\beta^2}{2}\int_0^1s\alpha(s)\dd s,
 \qquad P_\beta=\inf_\mu\cP_\beta(\mu).
\end{equation}
The infimum is attained by compactness and continuity, as proved in
Proposition~\ref{static-prop:scalar-minimizer}. These scalar approximation
and compactness arguments are independent of the pressure comparison;
the lower comparison below will use only finitely supported trial measures. Apply
Proposition~\ref{static-prop:matrix-bound} with $d=1$,
$D=Q_{\rm end}=1$, and $\lambda=0$. First for step measures and
then by approximation, this gives the upper comparison
\begin{equation}\label{static-eq:scalar-upper-bound}
 p_n(\beta):=\frac1n\E\log Z_{n,\beta}\le P_\beta.
\end{equation}

\begin{theorem}[Quantitative comparison]\label{static-thm:finite-size}
For every fixed $\beta>0$, there is $C_\beta<\infty$ such that for all
sufficiently large $n$,
\begin{equation}\label{static-eq:finite-size}
 P_\beta-C_\beta n^{-1/24}\le p_n(\beta)\le P_\beta.
\end{equation}
\end{theorem}

The remaining task is the lower comparison. We will choose an
interpolation path from expectations of overlaps, using a hierarchical
backward system. Such expectation-dependent choices are also a feature
of the adaptive interpolation method of Barbier and
Macris~\cite{BarbierMacris19}. Here we first prove an explicit overlap
error estimate for the hierarchy, without assuming a prior description
of the Gibbs measure, and then use it along the chosen path.
Mourrat's enriched-free-energy and Hamilton--Jacobi viewpoint
\cite[Section~3, equations~(3.24)--(3.27)]{Mourrat2020} provides related
motivation for organizing errors by hierarchical overlaps. This is a
methodological connection only: the direct covariance and concentration
argument and the quantitative bound here do not invoke a Hamilton--Jacobi
theorem.

\subsection{Concentration along a hierarchy}

To obtain an error estimate that can be summed over the levels, we
use equally spaced masses for the remainder of the section:
\[
 w_j=\frac1{k+1},\qquad x_j=\frac{j+1}{k+1}\quad(0\le j\le k).
\]
Fix $0\le t\le1$ and levels $0=q_{-1}<q_0<\cdots<q_k$.
Set $\Delta_j=\beta^2(q_j-q_{j-1})$ and use the terminal exponent
\begin{equation}\label{static-eq:adaptive-exponent}
 U(\sigma)=\sqrt t\,\beta\widehat H_n(\sigma)
              +\sum_{j=0}^k\sqrt{\Delta_j}\,z_j\cdot\sigma,
\end{equation}
where $z_0,\ldots,z_k$ are independent standard Gaussian vectors in
$\R^n$. Let $\mathcal A_j$ contain the prefix through level $j$,
including the quenched Hamiltonian disorder. Under the path law, put
\begin{equation}\label{static-eq:hierarchy-means}
 m_j=\E[\sigma\mid\mathcal A_j],\qquad
 M_j=\frac{\|m_j\|^2}{n},\qquad r_j=\E M_j.
\end{equation}
The conditional means $m_j$ form a martingale. Consequently
$0\le r_0\le\cdots\le r_k\le1$. They can also be obtained by
differentiating $F_j$ with respect to an additional deterministic field
coupled linearly to the spins. We begin by controlling the Gaussian
field energy under the same path law.

\begin{lemma}[Field-energy concentration]\label{static-lem:field-energy}
For every $0\le j\le k$, under the single-path law,
\begin{equation}\label{static-eq:field-energy-concentration}
 \Var\left(\frac{z_j\cdot\sigma}{n}\right)\le\frac1n.
\end{equation}
\end{lemma}

\begin{proof}
Condition first on the terminal spin $\sigma$. The product of the
tilts in Equation~\eqref{static-eq:hierarchy-tilt} and the terminal
Gibbs density telescopes. Thus the conditional joint density of all
standard Gaussian disorder variables and fields is proportional to
their original Gaussian density multiplied by
\begin{equation}\label{static-eq:fixed-spin-density}
 \exp\left\{U(\sigma)-\sum_{l=0}^kw_lF_l\right\}.
\end{equation}
Each $F_l$ is jointly convex in its prefix Gaussian variables. At the
terminal level this follows from log-sum-exp with linearly entering
variables. Backwards integration preserves convexity: expectation does
so at mass zero, and H\"older's inequality does so for a positive-mass
logarithmic expectation. Extend $F_l$ to all variables by making it
constant in its descendants; it is still convex. The weights $w_l$
are nonnegative, and $U(\sigma)$ is linear for fixed $\sigma$.
The negative logarithm of the conditional density therefore has Hessian
at least $I$. The linear specialization of
Lemma~\ref{static-lem:gaussian-tools} bounds the conditional variance
of $z_j\cdot\sigma/n$ by $1/n$.

To remove the conditioning, we must also check the conditional mean.
It is independent of $\sigma$, as follows from the gauge transformation
with signs $\eps_i=\sigma_i$:
\[
 J_{ab}\longmapsto\eps_a\eps_bJ_{ab}\quad(a<b),
 \qquad z_{l,i}\longmapsto\eps_i z_{l,i}\quad(0\le l\le k).
\]
Leave the diagonal variables fixed. This transformation preserves the
Gaussian law and carries the conditional law at $\sigma$ to the
conditional law at the all-plus spin. It also preserves
$z_j\cdot\sigma$. Thus the conditional means agree, and the law of
total variance proves Equation~\eqref{static-eq:field-energy-concentration}.
\end{proof}

\begin{lemma}[Average overlap concentration]\label{static-lem:average-overlap}
Let $\Delta_{\min}=\min_j\Delta_j$.  If $k\ge1$ and
$k/\sqrt{n\Delta_{\min}}\le1$, then
\begin{equation}\label{static-eq:average-overlap-concentration}
 \sum_{j=0}^kw_j\E_j^{(2)}(R-r_j)^2
 \le C\left(\frac{k}{\sqrt{n\Delta_{\min}}}+k^{-1/2}\right),
\end{equation}
where $C$ is an absolute constant.
\end{lemma}

\begin{proof}
The law of total variance splits the overlap error into two terms:
the fluctuation of its conditional mean $M_j$, and the variance left
after the prefix is fixed. We estimate these terms in that order.

\par\medskip\noindent\textbf{Fluctuation of the conditional mean.}\enspace
Fix $j<k$. Sample three paths $a,b,c$ with common prefix through
$j+1$ and conditionally independent continuations. If
$B=B(\sigma^a,\sigma^b,\sigma^c)$ satisfies $|B|\le1$, integration
by parts in the two adjacent Gaussian fields gives
\begin{equation}\label{static-eq:tree-ibp}
 \E B\left(\frac{z_j}{\sqrt{\Delta_j}}
           -\frac{z_{j+1}}{\sqrt{\Delta_{j+1}}}\right)
                    \cdot\frac{\sigma^a}{n}
 =-w_j\E B\frac{m_j\cdot\sigma^a}{n}.
\end{equation}
Here is the density calculation, including the common prefix. Relative
to the product Gaussian law, multiply $W_l$ along each common edge
$1\le l\le j+1$, along each of the three distinct descendant
branches, and then by the three terminal Gibbs probabilities. The
logarithm of this density is
\begin{align*}
 &\sum_{l=1}^{j+1}x_{l-1}(F_l-F_{l-1})\\
 &\quad+\sum_{v\in\{a,b,c\}}\left[
       \sum_{l=j+2}^kx_{l-1}(F_l^v-F_{l-1}^v)
                         +U^v(\sigma^v)-F_k^v\right].
\end{align*}
Apply the derivative difference
$\Delta_j^{-1/2}\partial_{z_{j,i}}
-\Delta_{j+1}^{-1/2}\partial_{z_{j+1,i}}$ coordinate by coordinate.
At level $j+1$ and below, every term depends on these fields only
through $\sqrt{\Delta_j}z_j+\sqrt{\Delta_{j+1}}z_{j+1}$, so its
two derivatives cancel. Terms above $j$ depend on neither field.
Only $F_j$ remains, with coefficient $x_{j-1}-x_j=-w_j$; the same
calculation applies at $j=0$ because $x_{-1}=0$. Its normalized
derivative is $m_{j,i}$. Throughout this differentiation the spin
labels and $B$ are fixed. Integration by parts and summation over
$i$ now prove Equation~\eqref{static-eq:tree-ibp}.

Let $Y$ be the quantity multiplying $B$ on the left-hand side of
Equation~\eqref{static-eq:tree-ibp}. Every branch has the single-path
marginal. Lemma~\ref{static-lem:field-energy} and the $L^2$ triangle
inequality therefore imply
\[
 \Var(Y)^{1/2}\le\frac{2}{\sqrt{n\Delta_{\min}}}.
\]
With $B=1$, Equation~\eqref{static-eq:tree-ibp} reads
$\E Y=-w_jr_j$. Subtract this mean, use $|B|\le1$, and recall
$w_j=1/(k+1)$ to obtain
\begin{equation}\label{static-eq:tree-centered}
 \left|\E B\left(\frac{m_j\cdot\sigma^a}{n}-r_j\right)\right|
 \le\frac{2(k+1)}{\sqrt{n\Delta_{\min}}}.
\end{equation}
Choose $B=R(\sigma^b,\sigma^c)$. Conditional on
$\mathcal A_{j+1}$ the three descendants are independent, so the
left-hand side becomes
\[
 \left|\E\left(\frac{m_j\cdot m_{j+1}}n-r_j\right)M_{j+1}\right|.
\]
We now replace $m_{j+1}$ by $m_j$ in both factors. The change in the
expectation is at most $3\sqrt{\gamma_j}$, where
\[
 \gamma_j=\frac1n\E\|m_{j+1}-m_j\|^2.
\]
This follows from the bounds $\|m_l\|\le\sqrt n$ and
$|M_j-r_j|\le1$. After replacement the expectation is
$\E(M_j-r_j)M_j=\Var(M_j)$. Martingale orthogonality gives
$\sum_{j<k}\gamma_j\le1$; averaging the errors and using
Cauchy--Schwarz proves
\begin{equation}\label{static-eq:mean-overlap-variance}
 \sum_{j<k}w_j\Var(M_j)
 \le C\left(\frac{k}{\sqrt{n\Delta_{\min}}}+k^{-1/2}\right).
\end{equation}

\par\medskip\noindent\textbf{Conditional variance.}\enspace
The conditional mean fluctuation is now controlled. To bound the
variance of the overlap of two descendants after their common prefix
is fixed, define
\[
 C_j=\Cov(\sigma\mid\mathcal A_j),\qquad
 D_j=\nabla^2F_j,\qquad
 V_j=\Cov(m_{j+1}\mid\mathcal A_j),
\]
where all conditional quantities are given $\mathcal A_j$, and
$D_j$ is the Hessian in an additive deterministic spin field. The
law of total covariance and the differentiated recursion give
\[
 C_j=\E[C_{j+1}\mid\mathcal A_j]+V_j,
 \qquad
 D_j=\E[D_{j+1}\mid\mathcal A_j]+x_jV_j.
\]
The terminal identity $D_k=C_k$, together with the nondecreasing
masses, allows backward induction and gives
\begin{equation}\label{static-eq:hessian-covariance-bound}
 D_j\succeq x_jC_j.
\end{equation}

For the next field, abbreviate $\Delta=\Delta_{j+1}$ and set
$y=\sqrt\Delta z_{j+1}$. Conditional on $\mathcal A_j$, its
density is proportional to
$\exp\{-\|y\|^2/(2\Delta)+x_jF_{j+1}\}$.
Integration by parts, using $\nabla_yF_{j+1}=m_{j+1}$, yields
\begin{align}
 \Cov(m_{j+1},y\mid\mathcal A_j)&=\Delta D_j,
 \label{static-eq:conditional-cross-covariance}\\
 \Cov(y\mid\mathcal A_j)&=\Delta I+\Delta^2x_jD_j.
 \label{static-eq:conditional-field-covariance}
\end{align}
For the first identity, use $\E[y\mid\mathcal A_j]=\Delta x_jm_j$
and subtract the product of means from the integration-by-parts
expression for $\E[m_{j+1}y^{\mathsf T}\mid\mathcal A_j]$.
A further application to $yy^{\mathsf T}$ proves the second identity.
The joint covariance matrix of $(m_{j+1},y)$ is positive semidefinite,
so its Schur complement gives
\begin{equation}\label{static-eq:covariance-schur-bound}
 V_j\succeq\Delta D_j(I+\Delta x_jD_j)^{-1}D_j.
\end{equation}

Order the eigenvalues of $C_j,D_j,V_j$ decreasingly and denote them
by $c_i,d_i,v_i$, respectively. Equation~\eqref{static-eq:hessian-covariance-bound}
implies $d_i\ge x_jc_i$. The function
$d\mapsto\Delta d^2/(1+\Delta x_jd)$ is increasing on $d\ge0$.
Apply the min--max principle to
Equation~\eqref{static-eq:covariance-schur-bound} to get
\[
 v_i\ge\frac{\Delta d_i^2}{1+\Delta x_jd_i}
      \ge\frac{\Delta x_j^2c_i^2}{1+\Delta x_j^2c_i}.
\]
At or above the threshold $c_i\ge(\Delta x_j^2)^{-1}$, the final
expression is at least $c_i/2$. Since also $c_i\le\tr C_j\le n$,
these eigenvalues contribute at most $2n\tr V_j$ to
$\sum_i c_i^2$. Those below the threshold contribute at most
$n/(\Delta x_j^2)$. Together the two bounds give
\begin{equation}\label{static-eq:conditional-square-trace}
 \frac{\tr C_j^2}{n^2}
 \le\frac{2\tr V_j}{n}
             +\frac1{n\Delta_{j+1}x_j^2}.
\end{equation}
Martingale orthogonality again supplies the summable quantity:
$\sum_{j<k}\E\tr V_j\le n$. Sum over $j$, using
$x_j=(j+1)/(k+1)$, to obtain
\begin{equation}\label{static-eq:average-square-trace}
 B_0:=\sum_{j<k}w_j\E\frac{\tr C_j^2}{n^2}
 \le\frac Ck+\frac{Ck}{n\Delta_{\min}}.
\end{equation}

For two independent descendants of the level-$j$ prefix, the
conditional overlap has mean $M_j$ and variance
\[
 \frac1{n^2}\left(\tr C_j^2+2m_j^{\mathsf T}C_jm_j\right).
\]
The contribution of the second term is at most
$2\sqrt{\tr C_j^2/n^2}$: use $\|m_j\|^2\le n$ and
$\|C_j\|_{\rm op}\le\sqrt{\tr C_j^2}$. Thus the averaged
conditional variance is at most $B_0+2\sqrt{B_0}$.
Under $k/\sqrt{n\Delta_{\min}}\le1$,
Equation~\eqref{static-eq:average-square-trace} gives $B_0\le C/k$.
Combine this estimate with the conditional-mean bound
\eqref{static-eq:mean-overlap-variance} through the law of total
variance. Finally, the level $k$ costs at most $4/(k+1)$ because
$|R-r_k|\le2$. These bounds prove
Equation~\eqref{static-eq:average-overlap-concentration}.
\end{proof}

\subsection{Adaptive interpolation and the finite-size error}

\begin{proof}[Proof of Theorem~\ref{static-thm:finite-size}]
The upper comparison was proved in
Equation~\eqref{static-eq:scalar-upper-bound}. To obtain the lower
comparison, we choose the field levels so that the pressure derivative
separates into a quadratic expression in the overlap means and the
overlap error controlled by Lemma~\ref{static-lem:average-overlap}.
After constructing that path, we will integrate the derivative and
identify its single-site endpoint with a scalar trial measure.

\par\medskip\noindent\textbf{Choosing the adaptive path.}\enspace
Consider provisionally a differentiable path $q(t)$ with positive gaps,
and set $f(t)=n^{-1}\E F_0(t,q(t))$. For a field gap in overlap units,
$a_j=q_j-q_{j-1}$, the hierarchy identity gives
\begin{equation}\label{static-eq:field-gap-derivative}
 \frac{\partial}{\partial a_j}\frac1n\E F_0
 =\frac{\beta^2}{2}
       \left(1-\sum_{l=j}^kw_lr_l\right).
\end{equation}
The Hamiltonian contribution to the derivative is
$\frac{\beta^2}{4}(1-\sum_lw_l\E_l^{(2)}R^2)$.
Summing the field contributions and using
$\sum_{j=0}^la_j'=q_l'$ therefore gives
\[
 f'(t)=\frac{\beta^2}{4}
       \left(1-\sum_{l=0}^kw_l\E_l^{(2)}R^2\right)
       +\frac{\beta^2}{2}
       \left(q_k'-\sum_{l=0}^kw_lr_lq_l'\right).
\]
This calculation suggests the system
\begin{equation}\label{static-eq:adaptive-ode}
 q_j'(t)=-r_j(t,q(t)).
\end{equation}
Indeed, since $\E_j^{(2)}R=r_j$, every positive-gap solution satisfies
\begin{equation}\label{static-eq:adaptive-pressure-derivative}
 f'(t)=\frac{\beta^2}{4}\left[
 1-2r_k+\sum_{j=0}^kw_jr_j^2
          -\sum_{j=0}^kw_j\E_j^{(2)}(R-r_j)^2\right].
\end{equation}
The final term is exactly the error already estimated; the remaining
terms can be bounded after integration using the differential equation.
We now construct a solution with gaps large enough to apply that estimate.
Choose
\[
 k=\lfloor n^{1/12}\rfloor,\qquad \eta=n^{-1/6},
 \qquad q_j(1)=(j+1)\eta\quad(0\le j\le k),
\]
and solve Equation~\eqref{static-eq:adaptive-ode} backwards from $t=1$.

\par\medskip\noindent\textbf{Existence and bounds for the adaptive path.}\enspace
We verify that this system has a solution throughout $[0,1]$ and
keeps every gap positive. On the open region of positive gaps, the
finite Gaussian recursions make $r_j$ continuous in $t$ and locally
Lipschitz in $q$. The following domination justifies those assertions.
Fix a compact subset of this region; in this local argument constants
may depend on $n,k,\beta$ and that subset. Uniformly in the prefix and
in $t\in[0,1]$, we have
$\|\nabla_{z_h}F_h\|=\sqrt{\Delta_h}\|m_h\|\le L$.
Subtract the value of $F_h$ at $z_h=0$ from both numerator and
denominator in the normalized tilt. This gives
\[
 W_h(z_h)\le C\exp(L\|z_h\|).
\]
The conditional first moments of subsequent fields are consequently
uniformly bounded. The derivative $\partial_{q_i}U$ is bounded by a
constant times $\|z_i\|+\|z_{i+1}\|$, with the second term omitted
at $i=k$. Differentiation through the recursions therefore yields
\[
 |\partial_{q_i}F_l|\le C\left(1+\sum_{h\le l}\|z_h\|\right).
\]
To apply this to $r_j$, represent it as the expectation of $R$ on
a two-path tree. Since terminal Gibbs probabilities are at most one,
the preceding estimates dominate its density and the $q_i$ derivative
of that density by $C(1+S)e^{LS}$, where $S$ is the sum of the norms
of all fields in this finite tree. The bound is Gaussian integrable,
uniformly in $t$ and in the chosen compact set of $q$ values. It
justifies the differentiated recursions and gives the required local
Lipschitz property. At $t=0$, continuity follows from the same bound
and continuity in $\sqrt t$; no derivative of $\sqrt t$ is needed.

Along the local solution, $r_j\ge r_{j-1}$ implies that each gap
$q_j-q_{j-1}$ is nondecreasing as time runs backwards, with
$r_{-1}=0$. Each gap therefore stays at least $\eta$. The bound
$0\le r_j\le1$ also prevents escape to infinity and gives
\begin{equation}\label{static-eq:adaptive-path-bounds}
 q_k(t)\le(k+1)\eta+1-t,
 \qquad\Delta_j(t)\ge\beta^2\eta.
\end{equation}
These estimates continue the local solution across the full interval.
They also verify the field-gap hypothesis of
Lemma~\ref{static-lem:average-overlap}. Thus, for all sufficiently
large $n$, uniformly along the path,
\begin{equation}\label{static-eq:adaptive-overlap-error}
 \sum_{j=0}^kw_j\E_j^{(2)}(R-r_j)^2
 \le C_\beta n^{-1/24}.
\end{equation}
Indeed, the two terms in that lemma are controlled by
$k/\sqrt{n\beta^2\eta}=O_\beta(n^{-1/3})$ and
$k^{-1/2}=O(n^{-1/24})$.

\par\medskip\noindent\textbf{Integration of the pressure derivative.}\enspace
The adaptive path has now been constructed and its overlap error is
bounded, so we can integrate
Equation~\eqref{static-eq:adaptive-pressure-derivative}.
Write $\delta=(k+1)\eta=O(n^{-1/12})$ and $Q_j=q_j(0)$.
The derivative in Equation~\eqref{static-eq:field-gap-derivative}
lies between $0$ and $\beta^2/2$. Consequently, removing the fields present at $t=1$
costs at most $\beta^2\delta/2$, and
\[
 f(1)\le p_n(\beta)+\frac{\beta^2}{2}\delta.
\]
The differential equation and Jensen's inequality also give
\[
 \int_0^1r_j(t)\dd t=Q_j-q_j(1),\qquad
 \int_0^1r_j(t)^2\dd t\ge\bigl(Q_j-q_j(1)\bigr)^2.
\]
The endpoint bounds are $0\le Q_j\le1+\delta$ and
$0\le q_j(1)\le\delta$. Integrate
Equation~\eqref{static-eq:adaptive-pressure-derivative}, apply these
bounds and the overlap error in
Equation~\eqref{static-eq:adaptive-overlap-error}, and conclude that
\begin{equation}\label{static-eq:adaptive-lower-bound}
 p_n(\beta)\ge
 f(0)+\frac{\beta^2}{4}
       \left(1-2Q_k+\sum_{j=0}^kw_jQ_j^2\right)
       -C_\beta n^{-1/24}.
\end{equation}
\par\medskip\noindent\textbf{Identification of a scalar trial measure.}\enspace
At $t=0$ there is no Hamiltonian coupling, so the recursion factors
over sites. Thus $f(0)$ is the scalar single-site value at levels
$Q_0,\ldots,Q_k$. The only possible obstruction to using these
levels as the support of a trial measure on $[0,1]$ is that some
may exceed one.

Replace each such level by $\overline Q_j=\min(Q_j,1)$.
Each gap can only decrease, and their total decrease is
$Q_k-\overline Q_k\le\delta$. Equation~\eqref{static-eq:field-gap-derivative}
bounds the change in the scalar value by $\beta^2\delta/2$.
The quadratic expression in
Equation~\eqref{static-eq:adaptive-lower-bound} changes by at most
$C_\beta\delta$ as well. We may therefore take $Q_k\le1$ within
the stated error, interpreting zero gaps by continuity.

Now define $\mu=\sum_{j=0}^kw_j\delta_{Q_j}$. Its distribution
function is zero up to $Q_0$, equals $x_{j-1}$ on each gap
$(Q_{j-1},Q_j)$, and equals one between $Q_k$ and $1$.
Over this final interval, the recursion maps $\log(2\cosh z)$ to
$\log(2\cosh z)+\beta^2(1-Q_k)/2$. Hence
\[
 \Phi_\alpha(0,0)=f(0)+\frac{\beta^2}{2}(1-Q_k).
\]
The penalty is identified by summation by parts:
\[
 \frac{\beta^2}{2}\int_0^1s\alpha(s)\dd s
 =\frac{\beta^2}{4}\left(1-\sum_{j=0}^kw_jQ_j^2\right).
\]
The expression preceding the error in
Equation~\eqref{static-eq:adaptive-lower-bound} is therefore exactly
$\cP_\beta(\mu)$. It is at least $P_\beta$, which completes the
lower comparison and the proof.
\end{proof}
